% META: {"id": "1NSI-RES-COURS-C3", "chapitre": "1NSI-RESEAUX", "type_objet": "cours", "capacites": ["P-ARCH-04A", "P-ARCH-04B"], "mode_creation": "ex_nihilo", "sources_inspiration": [], "status": "verified"}
\section{Capteurs, actionneurs et IHM programmée}

\definition[D4]{
  Un \textbf{capteur} mesure une grandeur physique (température, luminosité, présence...)
  et la convertit en une donnée numérique exploitable par un programme. Un
  \textbf{actionneur} reçoit une commande d'un programme et agit sur le monde physique
  (allumer une lampe, ouvrir un volet, déclencher un moteur...).
}

\exemple{Un thermostat connecté : le \textbf{capteur} de température mesure la
température ambiante ; le programme décide, selon cette mesure, d'activer ou non le
chauffage ; l'\textbf{actionneur} (relais électrique) coupe ou active effectivement le
chauffage.}

\propriete[Cahier des charges d'une IHM]{
  Réaliser une IHM (interface homme-machine) répondant à un \textbf{cahier des charges}
  consiste à traduire des règles énoncées en langage courant (« si la température descend
  sous $18\,^\circ\text{C}$, le chauffage s'active ») en un programme précis, testé sur des
  cas concrets.
}

\begin{codereference}{Thermostat programmé : capteur, décision, actionneur}
\begin{python}
etat = {"chauffage": False, "override_manuel": None}

def decider_chauffage(temperature, seuil=18):
    """Cahier des charges : le chauffage s'active si temperature < seuil,
    sauf si un override manuel est actif (bouton de la telecommande)."""
    if etat["override_manuel"] is not None:
        return etat["override_manuel"]
    return temperature < seuil

def actionner_chauffage(temperature, seuil=18):
    """Lit le capteur (simule par le parametre temperature), decide,
    et pilote l'actionneur (simule par etat["chauffage"])."""
    etat["chauffage"] = decider_chauffage(temperature, seuil)
    return etat["chauffage"]

def bouton_override(valeur):
    """Gestionnaire d'evenement du bouton manuel de la telecommande."""
    etat["override_manuel"] = valeur

print(actionner_chauffage(15))  # True : il fait froid, chauffage active
print(actionner_chauffage(22))  # False : il fait assez chaud
\end{python}
\end{codereference}

% BEGIN-VERIFY
% etat = {"chauffage": False, "override_manuel": None}
% def decider_chauffage(temperature, seuil=18):
%     if etat["override_manuel"] is not None:
%         return etat["override_manuel"]
%     return temperature < seuil
% def actionner_chauffage(temperature, seuil=18):
%     etat["chauffage"] = decider_chauffage(temperature, seuil)
%     return etat["chauffage"]
% def bouton_override(valeur):
%     etat["override_manuel"] = valeur
% assert actionner_chauffage(15) is True
% assert actionner_chauffage(22) is False
% END-VERIFY

\exemple{Le cahier des charges précise aussi qu'un utilisateur peut forcer manuellement le
chauffage (bouton \og override \fg{}), quelle que soit la température :}

% BEGIN-VERIFY
% etat = {"chauffage": False, "override_manuel": None}
% def decider_chauffage(temperature, seuil=18):
%     if etat["override_manuel"] is not None:
%         return etat["override_manuel"]
%     return temperature < seuil
% def actionner_chauffage(temperature, seuil=18):
%     etat["chauffage"] = decider_chauffage(temperature, seuil)
%     return etat["chauffage"]
% def bouton_override(valeur):
%     etat["override_manuel"] = valeur
% bouton_override(True)
% assert actionner_chauffage(25) is True  # force malgre une temperature elevee
% bouton_override(None)
% assert actionner_chauffage(25) is False  # retour au comportement automatique
% END-VERIFY

\erreurFrequente{Programmer une IHM en oubliant de tester les \textbf{cas limites} du
cahier des charges (température exactement égale au seuil, override activé puis désactivé).
Un cahier des charges mal testé peut sembler fonctionner sur l'exemple de démonstration tout
en étant incorrect dans des situations réelles fréquentes.}

\alamachine{Modifie le thermostat pour ajouter une \textbf{hystérésis} : le chauffage
s'active si la température descend sous $18\,^\circ\text{C}$, mais ne se désactive que
lorsqu'elle dépasse $20\,^\circ\text{C}$ (pour éviter des allumages/extinctions trop
fréquents). Teste ta fonction sur une séquence de températures
\lstinline{[15, 17, 19, 21, 19]}.}
